\paragraph{Observações (variante iii).}
\begin{itemize}
  \item \textbf{Infiltração:} com $I_r \le k_s$ toda a chuva entra até o fim (f10 com 36~mm/h, f50 com 9~mm/h, f100
    com 3,6~mm/h), e no f10 com 90~mm/h ($0{,}25k_s$) entram 93\,\%. Acima de $k_s$ a superfície alaga e a infiltração
    cai: 86\,\% no f50 com 36~mm/h ($= k_s$, alaga a partir de 11~h), 42\,\% no f50 com 90~mm/h (toda a superfície
    alagada em 4~h), 95\,\% no f100 com 9~mm/h e 29\,\% no f100 com 36~mm/h. No f100 com $2{,}5k_s$ quase tudo ainda
    entra porque o solo seco absorve acima de $k_s$ no começo (sorção).
  \item \textbf{Depois da chuva} a água infiltrada líquida diminui devagar (sai pela face de percolação) e os nós
    alagados voltam a 0--4.
  \item \textbf{Lençol freático} (Tabela~3 do artigo): em 24~h ele chega à superfície da crista no f50 com 90~mm/h
    (cota 29,9 em $x = 30$~m, superfície em 30) e quase no f10 com 180~mm/h (29,1): categoria I. No f50 com 36~mm/h ele
    sobe 6~m e aflora na face (cota 24,9 no meio da face, superfície em 25): categoria II. No f100 ele quase não sobe
    (0,2~m): a água fica numa zona rasa sobre a sucção (categoria IV). No f10 com 36~mm/h sobe 2,7~m.
  \item \textbf{Mecanismo:} a banda de ruptura do SRM fica a 4,0--4,6~m de profundidade nos casos perto da ruptura e
    a 5,2--6,5~m nas chuvas fracas; depois da chuva ela aprofunda (5,6--7,7~m em 48~h).
  \item \textbf{Numérica:} o transiente usou sempre o passo de 0,1~h (480 passos em 48~h, sem bissecção) com 4,0 a
    5,0 iterações de Newton por passo, e levou 114--131~s; o tempo total (1800--3000~s com 18 processos no PC) é
    quase todo dos 49 SRM.
\end{itemize}
